\section{Model}
\label{sec:model}

\para{Effects, outcomes, lag.}
A workflow runs effects $e_1,\dots,e_k$. For effect $e$, let $t_0$ be the time the ledger records \textsf{ISSUED}, immediately before the commit tap,
and let $t_{\mathrm{vis}}\in[t_0,\infty]$ be the time the effect becomes visible to end-to-end probes (in flight, server, rendering).
The hidden outcome is $X=\mathbf 1\{t_{\mathrm{vis}}<\infty\}$: $X{=}1$ means \emph{committed}, possibly later; $X{=}0$ means absent.
An effect cannot become visible while its app is dead or the network is down, so the lag is measured from an anchor $r_e$ that depends on the
effect's observable fault class $c$. For agent death $r_e=t_0$. For app death $r_e$ is the time the controller confirms its intent to relaunch the app. For connectivity loss
$r_e$ is the time the connection is restored. For a combined fault $r_e$ is the later of the last two. The lag of class $c$ is $T_c=\max(0,\,t_{\mathrm{vis}}-r_e)$ and its tail is
$G_c(w)=\PP(T_c>w\mid X{=}1,\text{class }c)$, with $G_c(w)\to0$ as $w\to\infty$ by definition of $X$. Calibration labels effects by the server instead, so an applied effect that no probe ever shows (K-9 with a Sent
folder, \S\ref{sec:eval-realsim}) enters the lag table as never visible. No wait then certifies, and the planner asks. An effect waits until the deadline $r_e+w_c$.
If the class is unknown, the planner uses the latest deadline over the classes. If an anchor it needs was not observed, it asks.
Waiting from $t_0$ instead waits too little after an outage. In measured-parameter simulation on Element with a 15\,s outage, product certification then retries
\CertClockTzeroDownFifteenProdPv\% of committed effects, against \CertClockDownFifteenProdPv\% when it waits from $r_e$ (Table~\ref{tab:cert-clock} in the appendix).
We write $T$ and $G$ for $T_c$ and $G_c$ when the class is fixed. Recovery ends with $C_e$ applied copies: a \emph{duplicate} is $C_e\ge2$
(retrying when $X{=}1$) and a \emph{loss} is $C_e=0$ (not retrying when $X{=}0$). Partial effects are modeled as several effects.

\para{Views and probe families.}
A \emph{view} is one rendering read by one reader (the history list through the accessibility tree, a detail screen through a VLM). Views of the same kind of evidence form a \emph{probe family}: hop-1 (the app's local outbox and sent list), nonce, digest
(a count of matching history items compared with a snapshot taken before the tap), and VLM. A read at time $t$ of view $v$ returns $Y\in\{0,1\}$ ($1$ = ``committed'').
We assume:
\begin{itemize}[leftmargin=*,topsep=1pt,itemsep=0pt]
\item[\textbf{A1}] (lag) A read of an end-to-end view at time $t$ has law $p^v_1$ if $X{=}1$ and $T\le t-r_e$, and $p^v_0$ otherwise:
an effect that is not yet visible looks absent. Hop-1 views observe the outbox entry, which exists from the tap on. A view qualifies as hop~1 only if its content
does not depend on rendering the end-to-end state. Element's send-state indicator does not qualify, because it is drawn with the lagged timeline row (\S\ref{sec:eval-real}).
Given $X$, the lag $T$ is independent of the readers' errors: a read depends on $T$ only through whether the effect is visible at $t$.
Element's \texttt{/sync} echo, which can show ``Sent'' before the app's own request returns, is an end-to-end view under A1. The echo is the event the server applied, so it appears only after the commit, and its arrival is when $T$ ends.
\item[\textbf{A2}] (independence across families) Given $X$ and visibility, readings of different families are independent.
Readings of views within a family may be arbitrarily dependent, and the planner reads each view at most once. Only product certification
(Proposition~\ref{prop:family}) uses A2, and joint certification (Proposition~\ref{prop:joint}) does not.
\end{itemize}
Re-reading a view adds nothing in this model (a reader tends to repeat its answer). For a binary view, $\epsilon_0=p_0(1)$ is the false-``committed'' rate and $\epsilon_1=p_1(0)$ the false-``absent'' rate.
We write $D_1=\KL(p_1\|p_0)$ and $D_0=\KL(p_0\|p_1)$, and the view is distinguishing if $D_0>0$. Theorem~\ref{thm:main} treats the classical case of a family whose
views are conditionally i.i.d., and \S\ref{sec:theory-corr} drops this.

\para{Actions, user, guarantees.}
After restart the planner may observe a view (cost $c_v$), retry, declare the effect done, ask the user, or compensate a
reversible effect. Asking is itself a noisy probe: the user answers wrongly with probability $\epsilon_u$ (higher when a look-alike is on screen), after a delay
$c_{\mathrm{ask}}$. A policy is a possibly randomized stopping rule plus a terminal action. It meets an $(\alpha,\beta)$ \emph{guarantee} if its own
decisions satisfy
\begin{equation}
\PP(\text{retry}\mid X{=}1)\le\alpha \quad\text{and}\quad \PP(\text{declare}\mid X{=}0)\le\beta .
\label{eq:guarantee}
\end{equation}
The guarantee is conditional on $X$, so it needs no prior. Exactly-once is $\alpha{=}\beta{=}0$. Errors introduced by wrong user answers are outside~\eqref{eq:guarantee}
and are reported separately (\S\ref{sec:eval}). \sys splits each effect's duplicate budget into a test share, a lag share (paid in waiting time) and a calibration share,
$\alpha=\alpha_{\mathrm{test}}+\alpha_{\mathrm{lag}}+\alpha_{\mathrm{cal}}$ with weights $0.5/0.4/0.1$, and $\beta=\beta_{\mathrm{test}}+\beta_{\mathrm{cal}}$ with $0.9/0.1$. These weights are
a design choice, and the guarantee holds for any split.

The end-to-end error adds what the user and the app do. Let $q_x$ be the probability that the planner asks about an effect with outcome $x$, and $\epsilon_x$ the user's
error rate on those asks. Let $r_x$ be the rate of errors that neither the planner nor the user causes, such as an app that retransmits without dedup. A duplicate or a loss
needs one of these three events, so a union bound gives
\begin{equation}
\begin{gathered}
\PP(\dup\mid X{=}1)\le\min(1,\alpha+q_1\epsilon_1+r_1),\\
\PP(\loss\mid X{=}0)\le\min(1,\beta+q_0\epsilon_0+r_0).
\end{gathered}
\label{eq:composed}
\end{equation}
A deployer who cares about end-to-end errors should read~\eqref{eq:composed}, not only the rate $P_{\mathrm v}$ at which the planner retries a committed effect by itself.
A mode that asks more meets~\eqref{eq:guarantee} more easily but pays $q_1\epsilon_1$. Our simulator assumes that no app retransmits without dedup, so $r_1{=}r_0{=}0$.
In measured-parameter simulation on Element, the duplicate bound is then \CertBoundElemProdDup\% for product and \CertBoundElemJointDup\% for joint (users erred on \CertBoundElemProdEpsOne\% and \CertBoundElemJointEpsOne\% of those asks: 5\%, or 10\% with a look-alike), and the loss bound is
\CertBoundElemProdLoss\% and \CertBoundElemJointLoss\%. Joint asks about \CertBoundElemJointQone\% of committed effects there (Table~\ref{tab:cert-bound} also gives the default world).
